\documentclass[11pt,a4paper]{article}
\usepackage[margin=1in]{geometry}
\usepackage{hyperref}
\usepackage{enumitem}
\usepackage{longtable}
\usepackage{booktabs}
\usepackage{xcolor}
\usepackage{titlesec}
\usepackage{fancyhdr}
\usepackage{tocloft}
\usepackage{listings}
\usepackage{verbatim}

\title{\textbf{Product Requirements Document (PRD)}\\
\textbf{AgeOfAgents}\\[0.5em]
\large The Ultimate Local-First RTS Command Center for AI Coding Agents \\
\large (Extended for Native macOS + Windows VM Fallback)}
\author{Dylan Kawalec (ChaoticBeautys) in collaboration with Grok AI}
\date{\today}

\hypersetup{
    colorlinks=true,
    linkcolor=blue,
    filecolor=magenta,
    urlcolor=cyan,
    pdftitle={AgeOfAgents PRD - Extended},
    pdfpagemode=FullScreen,
}

\begin{document}

\maketitle

\tableofcontents
\newpage

\section{Executive Summary \& Vision (Updated)}

\textbf{AgeOfAgents} is a complete fork and re-imagining of the open-source \textbf{openage} engine (your fork: \href{https://github.com/DylanCkawalec/openage}{https://github.com/DylanCkawalec/openage}) that transforms the classic \textit{Age of Empires II} real-time strategy (RTS) gameplay into the ultimate visual command center for orchestrating fleets of local AI coding agents.

\textbf{Key Update for macOS Users}: While the original openage documentation previously suggested limited macOS support, the current fork and upstream explicitly support **native macOS builds from source**. No Windows VM is required for core functionality. We will prioritize a seamless native macOS experience using Homebrew + Qt6. A Windows VM (via Parallels, UTM, or VMware) is documented only as a fallback for users who encounter rare graphics/driver issues on Apple Silicon.

The battlefield remains unchanged:
\begin{itemize}
    \item The **map = your entire massive monorepo + hundreds of internal repositories**.
    \item **Buildings = code modules, services, folders, and sub-repos**.
    \item **Units = AI Agents** (Coder, Planner, Reviewer, Executor, etc.).
    \item Resources, tech tree, and fog of war map directly to your protocol.
\end{itemize}

All raw ideas from AgentCraft (video: \href{https://www.youtube.com/watch?v=kR64LOqBBCU}{https://www.youtube.com/watch?v=kR64LOqBBCU}), Arcane Agents, VibeCraft, and GastownRTS are merged and replaced with production-grade local AI orchestration.

\textbf{Project Name}: \textbf{AgeOfAgents}

\textbf{Base Repository}: \href{https://github.com/DylanCkawalec/openage}{https://github.com/DylanCkawalec/openage} (your fork of SFTtech/openage).

\textbf{License}: GPLv3.

\textbf{Primary Target Platform}: macOS (native build) with full Windows/Linux parity.

\newpage
\section{Project Inspirations \& References (Unchanged)}

(Unchanged from previous PRD – see original for full list of links.)

\section{Core Concept: Creative Replacements (Unchanged)}

(Unchanged from previous PRD.)

\section{Functional Requirements (Unchanged)}

(Unchanged from previous PRD.)

\section{Non-Functional Requirements (Updated)}

\begin{itemize}
    \item Performance: Handle 100+ simultaneous agents on consumer hardware (including Apple Silicon M-series).
    \item Extensibility: Full Python API + nyan modding system.
    \item \textbf{Platform Support}: 
        \begin{itemize}
            \item \textbf{Native macOS} (primary – fully supported via source build).
            \item Windows (pre-built installers available upstream; VM fallback only).
            \item Linux (native).
        \end{itemize}
    \item Security: All execution inside user-controlled containers.
    \item \textbf{macOS-Specific}: GPU-accelerated OpenGL via Qt6, Apple Silicon optimized, no VM required for 99\% of users.
    \item Accessibility \& Usability: Hotkey-first + optional voice commands; seamless integration with macOS (menu bar, Dock, file sharing).
\end{itemize}

\newpage
\section{Technical Architecture Overview (Updated)}

\begin{itemize}
    \item \textbf{Core Engine}: Your fork of openage (C++ simulation + OpenGL/Qt6 UI).
    \item \textbf{New Backend Layer}: Cross-platform Node.js / Python bridge (Arcane-style tmux on macOS/Linux; Windows equivalent via Windows Subsystem for Linux or PowerShell).
    \item \textbf{Ollama Integration}: YAML-defined runtimes (native Ollama on macOS).
    \item \textbf{Modding}: nyan system for new agent roles.
    \item \textbf{State Persistence}: SQLite + openage savegame format.
    \item \textbf{Graphics}: Qt6 + OpenGL (macOS uses Metal translation layer automatically via system drivers).
\end{itemize}

\section{Platform Support \& macOS-Specific Setup (NEW SECTION)}

This section ensures AgeOfAgents runs **completely locally on your Mac** without mandatory Windows VM usage.

\subsection{1. Native macOS Build (Recommended \& Primary Path)}
openage (and therefore AgeOfAgents) supports native building on macOS. No pre-built packages exist, but the build process is fully documented and works on Apple Silicon and Intel Macs.

\textbf{Required Packages \& Drivers (Exact List from official macOS build guide)}:
\begin{itemize}
    \item Xcode Command Line Tools (Xcode \texttt{>= 12})
    \item Homebrew (latest)
    \item \texttt{brew update-reset \&\& brew update}
    \item \texttt{brew install --cask font-dejavu}
    \item \texttt{brew install cmake python3 libepoxy freetype fontconfig harfbuzz opus opusfile qt6 libogg libpng toml11 eigen llvm}
    \item Python packages (via pip):
      \texttt{pip3 install --upgrade --break-system-packages cython numpy mako lz4 pillow pygments setuptools toml}
    \item \textbf{nyan} (openage’s data description language) – build from \href{https://github.com/SFTtech/nyan}{https://github.com/SFTtech/nyan} using its own macOS instructions.
    \item \textbf{Graphics Drivers}: No extra install needed – macOS provides native OpenGL/Metal support. Qt6 automatically uses the system Metal layer for best performance on Apple Silicon.
    \item Optional but recommended for smooth UI: Latest macOS (Ventura or newer) + ensure OpenGL 4.1+ is available (standard on all modern Macs).
\end{itemize}

\textbf{Build Steps (AgeOfAgents-specific)}:
\begin{enumerate}
    \item Clone your fork: \texttt{git clone https://github.com/DylanCkawalec/openage.git \&\& cd openage}
    \item Follow \href{https://github.com/DylanCkawalec/openage/blob/master/doc/build_instructions/macos.md}{doc/build\_instructions/macos.md} exactly.
    \item Run \texttt{./configure --download-nyan --enable-debug} (or release).
    \item \texttt{make -j\$(sysctl -n hw.ncpu)}
    \item Launch: \texttt{./openage}
\end{enumerate}

\subsection{2. Docker-Based Build (Easiest for macOS)}
Use the official Docker build (works perfectly on macOS with Docker Desktop):
\begin{itemize}
    \item See \href{https://github.com/DylanCkawalec/openage/blob/master/doc/build_instructions/docker.md}{doc/build\_instructions/docker.md}.
    \item Mount your massive monorepo via Docker volumes for live access.
\end{itemize}

\subsection{3. Windows VM Sandbox (Fallback Only)}
If you encounter any rare Qt6/OpenGL rendering issues on macOS (e.g., specific Apple Silicon driver quirks), run inside a Windows VM:
\begin{itemize}
    \item Recommended: Parallels Desktop or UTM (free) with GPU passthrough enabled.
    \item Share your monorepo folder via Parallels Shared Folders or Samba.
    \item Use upstream Windows installer from \href{https://github.com/SFTtech/openage/releases}{SFTtech releases}.
    \item Additional Windows packages: Visual Studio 2022, vcpkg (for Qt6, etc.) – follow upstream Windows build guide.
\end{itemize}

\textbf{Why Native macOS is Preferred}: Full GPU acceleration, native file-system access to your monorepo, zero latency, and direct Ollama integration (Ollama runs natively on macOS with Apple Silicon acceleration).

\section{Implementation Roadmap (Updated with macOS Tasks)}

\subsection{Phase 0 – Fork \& Setup (1 week, macOS-focused)}
\begin{enumerate}
    \item Use your existing fork \href{https://github.com/DylanCkawalec/openage}{https://github.com/DylanCkawalec/openage}.
    \item Verify native macOS build succeeds with the exact dependency list above.
    \item Add macOS-specific CI workflow (.github/workflows/macos.yml) using GitHub Actions (already partially present upstream).
    \item Create AgeOfAgents modpack directory with macOS README.md containing the dependency table.
\end{enumerate}

(Phases 1–4 remain unchanged but now include macOS smoke tests in every phase.)

\section{Acceptance Criteria (Updated)}

The final system must:
\begin{itemize}
    \item Build and run **natively on macOS** using only the listed Homebrew/Qt6 packages (no VM required for standard use).
    \item Provide clear, copy-paste build instructions for macOS users.
    \item Support full monorepo access and Ollama integration on Apple Silicon.
    \item Fall back gracefully to Windows VM with shared-folder support.
    \item Preserve every idea from the four source projects while feeling like a polished, production-ready masterpiece.
\end{itemize}

\newpage
\section{Appendix A: Consolidated Feature DNA (Unchanged)}

(Refer to previous PRD.)

\textbf{This PRD is now fully extended and ready for implementation on your macOS machine.}

Let's build the future of AI coding together — starting with a native macOS AgeOfAgents that feels like it was made for your Mac.

\end{document}